% ==============================================================================
% Title Page
% ==============================================================================
\titleframe

\begin{frame}{Lecture Outline}
    \tableofcontents
\end{frame}

% ==============================================================================
% SECTION 1: MOTIVATION FOR TRANSFORMS
% ==============================================================================
\section{Motivation: Why Transform in Power Systems?}

\begin{frame}[allowframebreaks]{The Physical Reality of Power Systems}
    Power systems are composed of three-phase physical networks:
    \begin{itemize}
        \item Generators with rotating magnetic fields.
        \item Transmission lines with mutual magnetic and capacitive coupling between phases.
        \item Transformers with coupled iron cores and diverse winding connections.
        \item Loads that can be single-phase, unbalanced, non-linear, or dynamic.
    \end{itemize}

    \begin{alertblock}{The 4 Fundamental Mathematical Challenges}
        \begin{enumerate}
            \item \textbf{Electromagnetic Phase Coupling:} Off-diagonal elements in $\mathbf{Z}_{abc}$ and $\mathbf{Y}_{abc}$ prevent independent per-phase analysis.
            \item \textbf{Time Variance:} Sinusoidal AC steady states require tracking time-varying quantities; in rotating machines, inductances depend on rotor position $\theta(t)$.
            \item \textbf{Asymmetries and Faults:} Short-circuits (single line-to-ground, line-to-line) destroy phase symmetry.
            \item \textbf{High Dimensionality:} Thousands of interior buses with zero power injection burden dynamic simulations.
        \end{enumerate}
    \end{alertblock}
\end{frame}

\begin{frame}{What Transformations Accomplish: Machines \& Inverters}
    \begin{columns}
        \begin{column}{0.48\textwidth}
            \begin{block}{From AC to DC: Park Transform ($dq0$)}
                \begin{itemize}
                    \item "Freezes" rotating vectors onto a synchronous reference frame.
                    \item AC quantities become constant DC quantities in steady state.
                    \item Enables standard linear PI control for grid converters and machines.
                \end{itemize}
            \end{block}
        \end{column}
        \begin{column}{0.48\textwidth}
            \begin{block}{From 3D to 2D: Clarke Transform ($\alpha\beta0$)}
                \begin{itemize}
                    \item Projects $abc$ into orthogonal stationary axes.
                    \item Foundation of space vector modulation (SVPWM) and instantaneous power theory ($p\text{-}q$).
                \end{itemize}
            \end{block}
        \end{column}
    \end{columns}
\end{frame}

\begin{frame}{What Transformations Accomplish: Asymmetries \& Networks}
    \begin{columns}
        \begin{column}{0.48\textwidth}
            \begin{block}{Decoupling Asymmetries: Fortescue ($012$)}
                \begin{itemize}
                    \item Decomposes unbalanced 3-phase systems into 3 decoupled balanced sets.
                    \item Diagonalizes symmetric cyclic impedance matrices.
                    \item Essential for unsymmetrical fault analysis and protective relaying.
                \end{itemize}
            \end{block}
        \end{column}
        \begin{column}{0.48\textwidth}
            \begin{block}{Dimensional Reduction: Kron Reduction}
                \begin{itemize}
                    \item Eliminates passive zero-injection load buses.
                    \item Schur complement preserves exact multiport terminal behavior for stability studies.
                \end{itemize}
            \end{block}
        \end{column}
    \end{columns}
\end{frame}

\begin{frame}[shrink=10]{Taxonomy of Power System Transforms}
    \begin{table}
        \centering
        \scriptsize
        \begin{tabularx}{\textwidth}{l p{2.2cm} p{2.5cm} X}
            \toprule
            \textbf{Transformation} & \textbf{Domain} & \textbf{Primary Operation} & \textbf{Key Power System Use Case} \\
            \midrule
            \textbf{Fortescue} & Phasor / Freq. & Change of complex basis & Unbalanced faults, sequence networks, protection relays (46, 50N/51N) \\
            \addlinespace
            \textbf{Clarke} & Time / Instantaneous & 3-phase to 2-axis stationary ($\alpha\beta0$) & Space Vector PWM, instantaneous power theory ($p\text{-}q$), fast unbalance detection \\
            \addlinespace
            \textbf{Park} & Time / Instantaneous & 2-axis rotation to synchronous frame ($dq0$) & Machine modeling, vector control of grid converters (VSC/HVDC), SRF-PLL \\
            \addlinespace
            \textbf{Kron Reduction} & Admittance Network & Multiport matrix elimination (Schur complement) & Transient stability dynamic simulation, Ward equivalents, bus reduction \\
            \addlinespace
            \textbf{Full-Cycle DFT} & Discrete Time to Phasor & Discrete harmonic filtering & Digital distance protection, PMU synchrophasors (IEEE C37.118) \\
            \bottomrule
        \end{tabularx}
        \caption{Landscape of common power system transformations.}
    \end{table}
\end{frame}


% ==============================================================================
% SECTION 2: SYMMETRICAL COMPONENTS (FORTESCUE)
% ==============================================================================
\section{Direct, Inverse, and Zero Sequence Decomposition}

\begin{frame}[allowframebreaks]{The Symmetrical Components (Fortescue, 1918)}
    Charles LeGeyt Fortescue published in 1918:
    \begin{quote}
        \large
        \textit{"Method of Symmetrical Co-ordinates Applied to the Solution of Polyphase Networks"}
    \end{quote}

    \textbf{Fundamental Theorem:} Any set of $N$ unbalanced phasors can be decomposed into $N$ symmetrical (balanced) sets of phasors of different phase sequences.
 \newpage 
    For a 3-phase system ($N=3$), any unbalanced voltages $\bar{V}_a, \bar{V}_b, \bar{V}_c$ or currents $\bar{I}_a, \bar{I}_b, \bar{I}_c$ are uniquely decomposed into:
    \begin{enumerate}
        \item \textbf{Positive (Direct) Sequence ($1$ or $+$):} Three phasors of equal magnitude, displaced by $120^\circ$, with sequence $a \to b \to c$.
        \item \textbf{Negative (Inverse) Sequence ($2$ or $-$):} Three phasors of equal magnitude, displaced by $120^\circ$, with reversed sequence $a \to c \to b$.
        \item \textbf{Zero Sequence ($0$):} Three identical phasors of equal magnitude and identical phase angle.
    \end{enumerate}
\end{frame}

\begin{frame}{The Complex Operator $a$}
    The complex phase-shift operator $a$ rotates a phasor by $120^\circ$ counter-clockwise:
    \begin{equation}
        a = e^{j\frac{2\pi}{3}} = -\frac{1}{2} + j\frac{\sqrt{3}}{2} = 1\angle 120^\circ
    \end{equation}

    \begin{columns}
        \begin{column}{0.55\textwidth}
            \textbf{Key Algebraic Properties:}
            \begin{itemize}
                \item $a^2 = e^{j\frac{4\pi}{3}} = -\frac{1}{2} - j\frac{\sqrt{3}}{2} = 1\angle -120^\circ$
                \item $a^3 = 1, \quad a^4 = a$
                \item $1 + a + a^2 = 0$
                \item $1 - a = \sqrt{3} e^{-j 30^\circ}$
                \item $1 - a^2 = \sqrt{3} e^{j 30^\circ}$
            \end{itemize}
        \end{column}
        \begin{column}{0.45\textwidth}
            \begin{center}
                \begin{tikzpicture}[scale=0.95]
                    \draw[->, thick, gray!60] (-1.4,0) -- (1.4,0) node[right] {\tiny $\text{Re}$};
                    \draw[->, thick, gray!60] (0,-1.4) -- (0,1.4) node[above] {\tiny $\text{Im}$};
                    \draw[dashed, gray!50] (0,0) circle (1);
                    \draw[->, very thick, belgianblue] (0,0) -- (1,0) node[below right] {\scriptsize $1$};
                    \draw[->, very thick, myRed] (0,0) -- (-0.5, 0.866) node[above left] {\scriptsize $a$};
                    \draw[->, very thick, orange!80!black] (0,0) -- (-0.5, -0.866) node[below left] {\scriptsize $a^2$};
                \end{tikzpicture}
            \end{center}
        \end{column}
    \end{columns}
\end{frame}

\begin{frame}[allowframebreaks]{Mathematical Definition: Synthesis and Analysis}
    Expressing phase phasors from sequence components (\textbf{Synthesis}):
    \begin{align}
        \bar{V}_a &= \bar{V}_0 + \bar{V}_1 + \bar{V}_2 \\
        \bar{V}_b &= \bar{V}_0 + a^2 \bar{V}_1 + a \bar{V}_2 \\
        \bar{V}_c &= \bar{V}_0 + a \bar{V}_1 + a^2 \bar{V}_2
    \end{align}

    In matrix form:
    \begin{equation}
        \underbrace{\begin{bmatrix} \bar{V}_a \\ \bar{V}_b \\ \bar{V}_c \end{bmatrix}}_{\mathbf{V}_{abc}}
        =
        \underbrace{\begin{bmatrix}
            1 & 1 & 1 \\
            1 & a^2 & a \\
            1 & a & a^2
        \end{bmatrix}}_{\mathbf{T}}
        \underbrace{\begin{bmatrix} \bar{V}_0 \\ \bar{V}_1 \\ \bar{V}_2 \end{bmatrix}}_{\mathbf{V}_{012}}
    \end{equation}

    Inverting matrix $\mathbf{T}$ yields the \textbf{Analysis} transformation:
    \begin{equation}
        \mathbf{V}_{012} = \mathbf{T}^{-1} \mathbf{V}_{abc}
        =
        \frac{1}{3}
        \begin{bmatrix}
            1 & 1 & 1 \\
            1 & a & a^2 \\
            1 & a^2 & a
        \end{bmatrix}
        \begin{bmatrix} \bar{V}_a \\ \bar{V}_b \\ \bar{V}_c \end{bmatrix}
    \end{equation}

    Explicitly:
    \begin{align}
        \bar{V}_0 &= \frac{1}{3}(\bar{V}_a + \bar{V}_b + \bar{V}_c) \quad \text{\bf (Average phasor / ground unbalance)} \\
        \bar{V}_1 &= \frac{1}{3}(\bar{V}_a + a \bar{V}_b + a^2 \bar{V}_c) \quad \text{\bf (Direct / Positive component)} \\
        \bar{V}_2 &= \frac{1}{3}(\bar{V}_a + a^2 \bar{V}_b + a \bar{V}_c) \quad \text{\bf (Inverse / Negative component)}
    \end{align}
\end{frame}

\begin{frame}[allowframebreaks]{Physical Meaning of the Three Sequences}
    \begin{block}{Positive (Direct) Sequence $\bar{V}_1, \bar{I}_1$}
        \begin{itemize}
            \item Normal operating condition of the grid.
            \item Generates a stator magnetic field rotating \textbf{forward} at synchronous speed $\omega$.
            \item Delivers useful electromagnetic torque in motors and generators.
        \end{itemize}
    \end{block}

    \begin{alertblock}{Negative (Inverse) Sequence $\bar{V}_2, \bar{I}_2$}
        \begin{itemize}
            \item Caused by load unbalance, single-phase faults, or untransposed lines.
            \item Generates a magnetic field rotating in the \textbf{reverse direction} ($-\omega$).
            \item Relative velocity with respect to the rotor is $2\omega$ (100 Hz / 120 Hz).
            \item Induces dangerous double-frequency eddy currents in rotor surface and damper windings $\to$ \alert{\textbf{severe thermal overheating}} and pulsating torque.
        \end{itemize}
    \end{alertblock}

    \begin{block}{Zero Sequence $\bar{V}_0, \bar{I}_0$}
        \begin{itemize}
            \item All 3 phase currents $I_0$ oscillate completely in phase.
            \item Neutral / earth current: $I_n = I_a + I_b + I_c = 3 I_0$.
            \item \textbf{Can only circulate if a neutral connection or ground return path exists!}
            \item Generates a pulsating, non-rotating magnetic field in the airgap.
        \end{itemize}
    \end{block}
\end{frame}

\begin{frame}[allowframebreaks]{Decoupling of Symmetrical Networks}
    Consider a fully balanced three-phase line or equipment with self impedance $Z_s$ and mutual impedance $Z_m$:
    \begin{equation}
        \mathbf{Z}_{abc} = \begin{bmatrix}
            Z_s & Z_m & Z_m \\
            Z_m & Z_s & Z_m \\
            Z_m & Z_m & Z_s
        \end{bmatrix}
    \end{equation}
    In phase coordinates, phase $a$ voltage drop depends on currents in phases $b$ and $c$!

    Applying the Fortescue transformation $\mathbf{Z}_{012} = \mathbf{T}^{-1} \mathbf{Z}_{abc} \mathbf{T}$:
    \begin{equation}
        \mathbf{Z}_{012} = \begin{bmatrix}
            Z_0 & 0 & 0 \\
            0 & Z_1 & 0 \\
            0 & 0 & Z_2
        \end{bmatrix}
    \end{equation}
    where:
    \begin{align}
        Z_1 &= Z_s - Z_m \quad \text{\bf (Positive-sequence impedance)} \\
        Z_2 &= Z_s - Z_m = Z_1 \quad \text{\bf (Negative-sequence impedance)} \\
        Z_0 &= Z_s + 2 Z_m \quad \text{\bf (Zero-sequence impedance)}
    \end{align}

    \begin{alertblock}{The Magic of Symmetrical Components}
        In any cyclically symmetric equipment, \textbf{positive, negative, and zero sequence networks are completely decoupled!}
        \begin{itemize}
            \item Positive-sequence currents produce \textbf{only} positive-sequence voltage drops.
            \item Negative-sequence currents produce \textbf{only} negative-sequence voltage drops.
            \item Zero-sequence currents produce \textbf{only} zero-sequence voltage drops.
        \end{itemize}
    \end{alertblock}
\end{frame}

\begin{frame}[allowframebreaks]{Sequence Impedances of Power System Equipment}
    \textbf{1. Transmission Lines:}
    \begin{itemize}
        \item Transposed lines are static: $Z_1 = Z_2$.
        \item Zero sequence requires current return through the earth and ground wires.
        \item Because earth has finite conductivity and the return loop is large (Carson's formula):
        \begin{equation}
            Z_0 \approx 2 \text{ to } 3.5 \times Z_1
        \end{equation}
    \end{itemize}

    \textbf{2. Synchronous Machines:}
    \begin{itemize}
        \item Stator is rotating relative to fields: $Z_1 \neq Z_2 \neq Z_0$!
        \item $Z_1$: dynamic reactance ($X_d''$, $X_d'$, $X_d$ depending on time horizon).
        \item $Z_2$: determined by damper windings seeing a $2\omega$ backward field ($X_2 \approx \frac{X_d'' + X_q''}{2}$).
        \item $Z_0$: very small leakage reactance ($X_0 < X_2 < X_1$).
    \end{itemize}

    \textbf{3. Transformers:}
    \begin{itemize}
        \item $Z_1 = Z_2 = Z_{sc}$ (short-circuit impedance).
        \item $Z_0$ depends heavily on the core type and \alert{\textbf{winding connections}}!
        \begin{itemize}
            \item Delta ($\Delta$) winding: circulates zero sequence internally, but blocks it from reaching the external network!
            \item Ungrounded Wye (Y): infinite zero-sequence impedance (open circuit).
            \item Grounded Wye ($\text{Y}_n$): allows zero-sequence flow to ground.
        \end{itemize}
    \end{itemize}
\end{frame}

\begin{frame}[allowframebreaks]{Example 1: Numerical Decomposition of a Voltage Sag}
    Suppose phase $a$ suffers a voltage sag due to a remote fault, while phases $b$ and $c$ remain unaffected:
    \begin{equation}
        \bar{V}_a = 115 \angle 0^\circ \text{ V}, \quad \bar{V}_b = 230 \angle -120^\circ \text{ V}, \quad \bar{V}_c = 230 \angle 120^\circ \text{ V}
    \end{equation}

    Let us calculate the symmetrical components using $\mathbf{V}_{012} = \mathbf{T}^{-1} \mathbf{V}_{abc}$:
    \begin{align}
        \bar{V}_0 &= \frac{1}{3}(115 + 230 e^{-j 120^\circ} + 230 e^{j 120^\circ}) = \frac{115 - 230}{3} = \mathbf{-38.33 \angle 0^\circ \text{ V}} \\
        \bar{V}_1 &= \frac{1}{3}(115 + 230 e^{0^\circ} + 230 e^{0^\circ}) = \frac{115 + 460}{3} = \mathbf{191.67 \angle 0^\circ \text{ V}} \\
        \bar{V}_2 &= \frac{1}{3}(115 + 230 e^{j 120^\circ} + 230 e^{-j 120^\circ}) = \frac{115 - 230}{3} = \mathbf{-38.33 \angle 0^\circ \text{ V}}
    \end{align}

    \textbf{Interpretation:}
    \begin{itemize}
        \item The positive sequence dropped from $230\text{ V}$ to $191.67\text{ V}$.
        \item A negative-sequence voltage of $38.33\text{ V}$ (20\% unbalance) appeared.
        \item A zero-sequence voltage of $38.33\text{ V}$ appeared, which can be measured across the neutral point!
    \end{itemize}
\end{frame}

\begin{frame}[allowframebreaks]{Example 2: Unsymmetrical Fault Analysis}
    Symmetrical components allow solving unsymmetrical faults by connecting the three decoupled sequence networks at the fault point!

    \begin{columns}
        \begin{column}{0.48\textwidth}
            \begin{block}{Single Line-to-Ground (SLG)}
                Phase $a$ to ground with fault impedance $Z_f$:
                \begin{equation}
                    I_b = 0, \quad I_c = 0, \quad V_a = Z_f I_a
                \end{equation}
                Transforming to sequence currents:
                \begin{equation}
                    I_0 = I_1 = I_2 = \frac{1}{3} I_a
                \end{equation}
                $\implies$ \alert{\textbf{Networks in SERIES!}}
                \begin{equation}
                    I_1 = \frac{E_a}{Z_1 + Z_2 + Z_0 + 3 Z_f}
                \end{equation}
            \end{block}
        \end{column}

        \begin{column}{0.48\textwidth}
            \begin{block}{Line-to-Line (L-L) Fault}
                Phases $b$ and $c$ shorted together:
                \begin{equation}
                    I_a = 0, \quad I_b = -I_c, \quad V_b = V_c
                \end{equation}
                Transforming to sequence currents:
                \begin{equation}
                    I_0 = 0, \quad I_1 = -I_2
                \end{equation}
                $\implies$ \alert{\textbf{Networks in PARALLEL!}}
                Zero sequence is disconnected.
                \begin{equation}
                    I_1 = \frac{E_a}{Z_1 + Z_2 + Z_f}
                \end{equation}
            \end{block}
        \end{column}
    \end{columns}
\end{frame}

\begin{frame}{Example 3: Protective Relaying Applications}
    Symmetrical components are directly used to configure numerical protection relays:
    \vfill
    \begin{columns}
        \begin{column}{0.48\textwidth}
            \begin{block}{Negative-Sequence Overcurrent (46)}
                \begin{itemize}
                    \item Monitors magnitude of $|I_2|$.
                    \item Protects generators against rotor overheating caused by unbalances.
                    \item Tripping characteristic:
                    \begin{equation}
                        I_2^2 t \le K
                    \end{equation}
                    where $K$ is the generator thermal capability constant.
                \end{itemize}
            \end{block}
        \end{column}

        \begin{column}{0.48\textwidth}
            \begin{block}{Ground Overcurrent (50N / 51N)}
                \begin{itemize}
                    \item Monitors residual current $3 I_0 = I_a + I_b + I_c$.
                    \item In balanced conditions, $3 I_0 \approx 0$.
                    \item Any ground fault immediately injects large $I_0$.
                    \item Detects high-impedance earth faults that standard overcurrent (50/51) misses completely.
                \end{itemize}
            \end{block}
        \end{column}
    \end{columns}
\end{frame}


% ==============================================================================
% SECTION 3: CLARKE TRANSFORM
% ==============================================================================
\section{\texorpdfstring{The Clarke Transform ($\alpha\beta0$ Frame)}{The Clarke Transform (alpha-beta-0 Frame)}}

\begin{frame}[allowframebreaks]{\texorpdfstring{The Clarke Transform (Edith Clarke, 1938)}{The Clarke Transform (Edith Clarke, 1938)}}
    Edith Clarke introduced the $\alpha\beta0$ transformation to simplify the calculation of unsymmetrical three-phase circuits.

    \textbf{Core Concept:}
    \begin{itemize}
        \item In a balanced three-phase system, $v_a(t) + v_b(t) + v_c(t) = 0$.
        \item Thus, three variables are redundant: there are only \textbf{2 degrees of freedom}!
        \item The Clarke transform converts the 3 stationary physical axes ($a, b, c$, separated by $120^\circ$) into 2 orthogonal stationary axes ($\alpha, \beta$, separated by $90^\circ$), plus a zero-sequence axis ($0$).
    \end{itemize}

    \begin{center}
        \begin{tikzpicture}[scale=1.1]
            \draw[->, thick, gray!60] (-1.8,0) -- (2.3,0) node[right] {\scriptsize $\alpha$ (aligned with $a$)};
            \draw[->, thick, gray!60] (0,-1.3) -- (0,1.8) node[above] {\scriptsize $\beta$ ($+90^\circ$)};
            \draw[->, very thick, belgianblue] (0,0) -- (1.8,0) node[below] {\scriptsize \bf Phase $a$};
            \draw[->, very thick, orange!80!black] (0,0) -- (-0.9, 1.56) node[above left] {\scriptsize \bf Phase $b$ ($+120^\circ$)};
            \draw[->, very thick, myRed] (0,0) -- (-0.9, -1.56) node[below left] {\scriptsize \bf Phase $c$ ($-120^\circ$)};
            \draw[->, very thick, purple] (0,0) -- (1.4, 0.9) node[above right] {\scriptsize $\vec{v}_{\alpha\beta}(t)$};
            \draw[dashed, purple] (1.4, 0.9) -- (1.4, 0) node[below] {\tiny $v_\alpha$};
            \draw[dashed, purple] (1.4, 0.9) -- (0, 0.9) node[left] {\tiny $v_\beta$};
        \end{tikzpicture}
    \end{center}
\end{frame}

\begin{frame}[allowframebreaks]{Mathematical Formulations: Amplitude vs Power Invariant}
    \begin{block}{Amplitude-Invariant Clarke Transform (Standard in Control)}
        Preserves the peak amplitude of phase quantities ($\hat{V}_\alpha = \hat{V}_a$ in balanced conditions):
        \begin{equation}
            \begin{bmatrix} v_\alpha \\ v_\beta \\ v_0 \end{bmatrix}
            =
            \frac{2}{3}
            \begin{bmatrix}
                1 & -\frac{1}{2} & -\frac{1}{2} \\
                0 & \frac{\sqrt{3}}{2} & -\frac{\sqrt{3}}{2} \\
                \frac{1}{2} & \frac{1}{2} & \frac{1}{2}
            \end{bmatrix}
            \begin{bmatrix} v_a \\ v_b \\ v_c \end{bmatrix}
        \end{equation}

        \textbf{Inverse Transform:}
        \begin{equation}
            \begin{bmatrix} v_a \\ v_b \\ v_c \end{bmatrix}
            =
            \begin{bmatrix}
                1 & 0 & 1 \\
                -\frac{1}{2} & \frac{\sqrt{3}}{2} & 1 \\
                -\frac{1}{2} & -\frac{\sqrt{3}}{2} & 1
            \end{bmatrix}
            \begin{bmatrix} v_\alpha \\ v_\beta \\ v_0 \end{bmatrix}
        \end{equation}
    \end{block}

    \begin{block}{Power-Invariant Clarke Transform (Concordia Transform)}
        Uses an orthonormal matrix ($\mathbf{C}^{-1} = \mathbf{C}^T$) to preserve instantaneous power $p(t) = v_a i_a + v_b i_b + v_c i_c = v_\alpha i_\alpha + v_\beta i_\beta + v_0 i_0$:
        \begin{equation}
            \mathbf{C}_{\text{power}}
            =
            \sqrt{\frac{2}{3}}
            \begin{bmatrix}
                1 & -\frac{1}{2} & -\frac{1}{2} \\
                0 & \frac{\sqrt{3}}{2} & -\frac{\sqrt{3}}{2} \\
                \frac{1}{\sqrt{2}} & \frac{1}{\sqrt{2}} & \frac{1}{\sqrt{2}}
            \end{bmatrix}
        \end{equation}
    \end{block}
\end{frame}

\begin{frame}[allowframebreaks]{The Space Vector Trajectory}
    The complex space vector in the stationary plane is:
    \begin{equation}
        \vec{v}_{\alpha\beta}(t) = v_\alpha(t) + j v_\beta(t)
    \end{equation}

    \begin{columns}
        \begin{column}{0.48\textwidth}
            \begin{block}{Balanced Conditions}
                \begin{itemize}
                    \item $v_a(t) = \hat{V} \cos(\omega t)$
                    \item $v_\alpha(t) = \hat{V} \cos(\omega t)$
                    \item $v_\beta(t) = \hat{V} \sin(\omega t)$
                    \item $v_0(t) = 0$
                \end{itemize}
                $\implies$ Traces a \alert{\textbf{perfect circle of radius $\hat{V}$}} rotating at angular speed $\omega$.
            \end{block}
        \end{column}

        \begin{column}{0.48\textwidth}
            \begin{alertblock}{Unbalanced Conditions}
                \begin{itemize}
                    \item Forward circle + backward circle:
                    \begin{equation*}
                        \vec{v}(t) = \hat{V}_1 e^{j(\omega t + \phi_1)} + \hat{V}_2 e^{-j(\omega t - \phi_2)}
                    \end{equation*}
                \end{itemize}
                $\implies$ Traces an \alert{\textbf{ellipse}} on the $\alpha\beta$ plane. The ellipticity directly measures voltage unbalance!
            \end{alertblock}
        \end{column}
    \end{columns}
\end{frame}

\begin{frame}[allowframebreaks]{Practical Example: Space Vector PWM (SVPWM)}
    In voltage-source converters (wind, solar, battery inverters, HVDC-VSC):
    \begin{itemize}
        \item A 2-level inverter has 8 switching states ($S_a, S_b, S_c \in \{0, 1\}$).
        \item In the Clarke $\alpha\beta$ plane, the 8 states map to \textbf{6 active voltage vectors} forming a regular hexagon + 2 zero vectors at the origin!
    \end{itemize}

    \begin{center}
        \begin{tikzpicture}[scale=0.95]
            \foreach \angle / \vecname in {0/$V_1(100)$, 60/$V_2(110)$, 120/$V_3(010)$, 180/$V_4(011)$, 240/$V_5(001)$, 300/$V_6(101)$} {
                \draw[->, thick, belgianblue] (0,0) -- (\angle:1.8) node[anchor=\angle+180] {\tiny \vecname};
            }
            \draw[dashed, gray] (0:1.8) -- (60:1.8) -- (120:1.8) -- (180:1.8) -- (240:1.8) -- (300:1.8) -- cycle;
            \draw[fill=myRed] (0,0) circle (0.05) node[below right] {\tiny $V_0, V_7$};
            \node at (30:1.0) {\tiny \bf Sector I};
            \node at (90:1.0) {\tiny \bf Sector II};
            \draw[->, very thick, purple] (0,0) -- (40:1.4) node[above] {\tiny $\vec{V}_{\text{ref}}$};
        \end{tikzpicture}
    \end{center}
    \textbf{Result:} Generating the target reference voltage $\vec{V}_{\text{ref}}$ simply requires duty-cycle calculation between adjacent sector boundary vectors.
\end{frame}


% ==============================================================================
% SECTION 4: PARK TRANSFORM
% ==============================================================================
\section{\texorpdfstring{The Park Transform ($dq0$ Frame) and Grid Synchronization}{The Park Transform (dq0 Frame) and Grid Synchronization}}

\begin{frame}[allowframebreaks]{The Park Transform (Robert H. Park, 1929)}
    In 1929, Robert H. Park published a landmark paper in power engineering:
    \begin{quote}
        \textit{"Two-reaction theory of synchronous machines"}
    \end{quote}

    \textbf{The Fundamental Machine Problem:}
    \begin{itemize}
        \item In a salient-pole synchronous generator, the air-gap magnetic permeance varies with rotor angle $\theta(t)$.
        \item Stator self and mutual inductances $L_{aa}(\theta), L_{ab}(\theta)$ are \textbf{periodic functions of rotor position}:
        \begin{equation}
            L_{aa}(\theta) = L_s + L_m \cos(2\theta)
        \end{equation}
        \item This makes machine differential equations time-varying, coupled, and non-linear!
    \end{itemize}

    \textbf{Park's Brilliant Solution:}
    Attach the reference frame to the \textbf{rotating rotor} (rotating at speed $\omega = \frac{d\theta}{dt}$).
    In this rotating $dq$ frame, the rotor appears stationary $\implies$ \alert{\textbf{all inductances become constant!}}
\end{frame}

\begin{frame}[allowframebreaks]{Mathematical Definition: Stationary to Rotating Frame}
    The Park transform rotates the stationary $\alpha\beta$ axes by angle $\theta(t)$:
    \begin{equation}
        \begin{bmatrix} v_d \\ v_q \end{bmatrix}
        =
        \begin{bmatrix}
            \cos\theta & \sin\theta \\
            -\sin\theta & \cos\theta
        \end{bmatrix}
        \begin{bmatrix} v_\alpha \\ v_\beta \end{bmatrix}
    \end{equation}

    \textbf{Inverse Park Transform:}
    \begin{equation}
        \begin{bmatrix} v_\alpha \\ v_\beta \end{bmatrix}
        =
        \begin{bmatrix}
            \cos\theta & -\sin\theta \\
            \sin\theta & \cos\theta
        \end{bmatrix}
        \begin{bmatrix} v_d \\ v_q \end{bmatrix}
    \end{equation}

    \begin{center}
        \begin{tikzpicture}[scale=1.1]
            \draw[->, thick, gray!60] (-0.5,0) -- (2.3,0) node[right] {\scriptsize $\alpha$};
            \draw[->, thick, gray!60] (0,-0.5) -- (0,2.0) node[above] {\scriptsize $\beta$};
            \draw[dashed, gray] (0,0) -- (30:2.3);
            \draw[->, thick, belgianblue] (0,0) -- (30:2.0) node[above right] {\scriptsize $d$ (direct)};
            \draw[->, thick, orange!80!black] (0,0) -- (120:2.0) node[above left] {\scriptsize $q$ (quadrature, $+90^\circ$)};
            \draw[->] (0.8,0) arc (0:30:0.8) node[midway, right] {\scriptsize $\theta(t)$};
        \end{tikzpicture}
    \end{center}
\end{frame}

\begin{frame}{Direct $abc \to dq0$ Formulation}
    Combining the Clarke and Park matrices yields the direct $abc \to dq0$ transformation:
    \begin{equation}
        \begin{bmatrix} v_d \\ v_q \\ v_0 \end{bmatrix}
        =
        \frac{2}{3}
        \begin{bmatrix}
            \cos\theta & \cos(\theta - \frac{2\pi}{3}) & \cos(\theta + \frac{2\pi}{3}) \\
            -\sin\theta & -\sin(\theta - \frac{2\pi}{3}) & -\sin(\theta + \frac{2\pi}{3}) \\
            \frac{1}{2} & \frac{1}{2} & \frac{1}{2}
        \end{bmatrix}
        \begin{bmatrix} v_a \\ v_b \\ v_c \end{bmatrix}
    \end{equation}
    where $\theta(t) = \int \omega(t) dt$.
\end{frame}

\begin{frame}{The Ultimate Magic: AC Becomes DC}
    For balanced three-phase sinusoids with phase angle $\phi$:
    \begin{align}
        v_a(t) &= \hat{V} \cos(\omega t + \phi) \\
        v_b(t) &= \hat{V} \cos(\omega t + \phi - 2\pi/3) \\
        v_c(t) &= \hat{V} \cos(\omega t + \phi + 2\pi/3)
    \end{align}
    Setting the rotating frame angle to $\theta(t) = \omega t$:
    \begin{alertblock}{Steady-State DC Quantities}
        \begin{equation}
            \mathbf{v_d = \hat{V} \cos\phi = \text{constant}}, \qquad \mathbf{v_q = \hat{V} \sin\phi = \text{constant}}, \qquad \mathbf{v_0 = 0}
        \end{equation}
    \end{alertblock}
    Sinusoids at $50\text{ Hz}$ or $60\text{ Hz}$ appear as completely flat DC levels to the controller!
\end{frame}

\begin{frame}[allowframebreaks]{Control Advantage: Decoupled Vector Control}
    In grid-connected inverters (solar, wind, battery, HVDC):
    \begin{itemize}
        \item Instantaneous active power $P$ and reactive power $Q$ in the $dq$ frame (amplitude-invariant):
        \begin{equation}
            P = \frac{3}{2}(v_d i_d + v_q i_q), \qquad Q = \frac{3}{2}(v_q i_d - v_d i_q)
        \end{equation}
        \item If we align the $d$-axis with the grid voltage vector such that $v_d = \hat{V}$ and $v_q = 0$:
        \begin{equation}
            \mathbf{P = \frac{3}{2} v_d i_d \propto i_d}, \qquad \mathbf{Q = -\frac{3}{2} v_d i_q \propto -i_q}
        \end{equation}
    \end{itemize}

    \begin{alertblock}{Total Decoupling}
        \begin{itemize}
            \item $i_d$ independently controls Active Power $P$ (or DC-link voltage).
            \item $i_q$ independently controls Reactive Power $Q$ (or grid AC voltage).
            \item Because $i_d, i_q$ are DC quantities in steady state, simple \textbf{Proportional-Integral (PI) controllers} achieve zero steady-state tracking error without phase lag!
        \end{itemize}
    \end{alertblock}
\end{frame}

\begin{frame}[allowframebreaks]{Behavior Under Grid Unbalance: The $2\omega$ Ripple}
    What happens if the grid voltages are unbalanced?
    Recall that unbalanced voltages contain both positive ($\bar{V}_1$) and negative ($\bar{V}_2$) sequences:
    \begin{align}
        \vec{v}_{\alpha\beta}(t) &= \hat{V}_1 e^{j(\omega t + \phi_1)} + \hat{V}_2 e^{-j(\omega t - \phi_2)}
    \end{align}

    Rotating with the forward angle $\theta = \omega t$:
    \begin{equation}
        v_{dq}(t) = \vec{v}_{\alpha\beta}(t) e^{-j\omega t} = \underbrace{\hat{V}_1 e^{j\phi_1}}_{\text{\bf DC component}} + \underbrace{\hat{V}_2 e^{-j(2\omega t - \phi_2)}}_{\text{\bf Double-frequency AC ripple at } \mathbf{2\omega}}
    \end{equation}

    \begin{alertblock}{Key Diagnostic Insight}
        \begin{itemize}
            \item Positive sequence appears as a pure \textbf{DC level} in the $dq$ frame.
            \item Negative sequence appears as a \textbf{100 Hz (or 120 Hz) sinusoidal oscillation} in $v_d$ and $v_q$!
            \item Dual-frame controllers (DDSRF-PLL) separate sequences to ride through asymmetric faults.
        \end{itemize}
    \end{alertblock}
\end{frame}

\begin{frame}[allowframebreaks]{Grid Synchronization: The SRF-PLL}
    To execute the Park transform, we must track the exact grid voltage angle $\theta(t)$ in real time.
    This is achieved by the \textbf{Synchronous Reference Frame Phase-Locked Loop (SRF-PLL)}:

    \begin{center}
        \begin{tikzpicture}[scale=0.9]
            \node[draw, rectangle, fill=headerbg, minimum height=0.8cm, text width=1.3cm, align=center] (park) at (0,0) {\tiny Park};
            \node[draw, circle, fill=white, inner sep=1pt] (sum) at (2.0,0) {\tiny $-$};
            \node[draw, rectangle, fill=headerbg, minimum height=0.8cm, text width=1.4cm, align=center] (pi) at (3.8,0) {\tiny PI Controller\\$k_p + \frac{k_i}{s}$};
            \node[draw, circle, fill=white, inner sep=1pt] (sumw) at (5.7,0) {\tiny $+$};
            \node[draw, rectangle, fill=headerbg, minimum height=0.8cm, text width=1.3cm, align=center] (int) at (7.3,0) {\tiny Integrator\\$\frac{1}{s}$};

            \draw[->, thick] (-1.4,0.15) -- (park.170) node[midway, above] {\tiny $v_\alpha$};
            \draw[->, thick] (-1.4,-0.15) -- (park.190) node[midway, below] {\tiny $v_\beta$};
            \draw[->, thick] (park.east) -- (sum.west) node[midway, above] {\tiny $v_q$};
            \draw[->, thick] (2.0, -0.6) -- (sum.south) node[left] {\tiny $v_q^* = 0$};
            \draw[->, thick] (sum.east) -- (pi.west);
            \draw[->, thick] (pi.east) -- (sumw.west) node[midway, above] {\tiny $\Delta\omega$};
            \draw[->, thick] (5.7, 0.6) -- (sumw.north) node[above] {\tiny $\omega_0$};
            \draw[->, thick] (sumw.east) -- (int.west) node[midway, above] {\tiny $\omega$};
            \draw[->, thick] (int.east) -- (8.4,0) node[right] {\tiny $\theta(t)$};
            \draw[->, thick] (8.0,0) -- (8.0,-0.9) -- (0,-0.9) -- (park.south);
        \end{tikzpicture}
    \end{center}

    \textbf{Operating Principle:} If the estimated angle $\theta$ matches the grid angle, $v_q = 0$. The PI controller adjusts estimated frequency $\omega$ until error $v_q \to 0$.
\end{frame}


% ==============================================================================
% SECTION 5: BRIDGES BETWEEN TRANSFORMS
% ==============================================================================
\section{Bridges and Connections Between Transforms}

\begin{frame}[shrink=10]{Comparing Fortescue, Clarke, and Park}
    \begin{table}
        \centering
        \scriptsize
        \begin{tabularx}{\textwidth}{l p{2.8cm} p{3.0cm} X}
            \toprule
            \textbf{Feature} & \textbf{Fortescue ($012$)} & \textbf{Clarke ($\alpha\beta0$)} & \textbf{Park ($dq0$)} \\
            \midrule
            \textbf{Domain} & Complex Phasor (Frequency) & Instantaneous Time & Instantaneous Time \\
            \addlinespace
            \textbf{Coordinates} & Complex scalars $\bar{V}_0, \bar{V}_1, \bar{V}_2$ & Stationary orthogonal $(\alpha, \beta, 0)$ & Synchronous rotating $(d, q, 0)$ \\
            \addlinespace
            \textbf{Balanced Steady State} & $\bar{V}_1 = V_{\text{rms}} \angle \phi$, $\bar{V}_2 = 0, \bar{V}_0 = 0$ & Rotating circle of radius $\hat{V}$ at $\omega$ & Constant DC values ($V_d, V_q = \text{const}$) \\
            \addlinespace
            \textbf{Unbalanced Steady State} & Non-zero $\bar{V}_2$ and $\bar{V}_0$ (constant phasors) & Ellipse on $\alpha\beta$ plane & DC value + $2\omega$ AC ripple in $d$ and $q$ \\
            \addlinespace
            \textbf{Primary Strengths} & Frequency-domain fault calculation, symmetrical networks, protection & Space Vector PWM, instant power calculation, converter modeling & Linear PI control, machine flux/torque decoupling, PLL synchronization \\
            \bottomrule
        \end{tabularx}
        \caption{Unified comparison of coordinate transformations.}
    \end{table}
\end{frame}

\begin{frame}[allowframebreaks]{Mapping Between Domains}
    Under fundamental frequency operation, exact mappings link these three representations:

    \begin{block}{Clarke $\longleftrightarrow$ Fortescue}
        The complex space vector relates directly to positive and negative sequence phasors:
        \begin{equation}
            \vec{v}_{\alpha\beta}(t) = \sqrt{2}\left(\bar{V}_1 e^{j\omega t} + \bar{V}_2^* e^{-j\omega t}\right)
        \end{equation}
        \begin{itemize}
            \item Forward rotating part $\to \bar{V}_1$
            \item Backward rotating part $\to \bar{V}_2$
            \item Homopolar part $v_0(t) = \sqrt{2} \operatorname{Re}\{\bar{V}_0 e^{j\omega t}\}$
        \end{itemize}
    \end{block}

    \begin{block}{Park $\longleftrightarrow$ Fortescue}
        Demodulating in the rotating frame with $\theta = \omega t$:
        \begin{equation}
            v_d(t) + j v_q(t) = \vec{v}_{\alpha\beta}(t) e^{-j\omega t} = \sqrt{2} \bar{V}_1 + \sqrt{2} \bar{V}_2^* e^{-j 2\omega t}
        \end{equation}
        \begin{itemize}
            \item Mean (DC) value gives \textbf{positive sequence}!
            \item Amplitude of $2\omega$ ripple gives \textbf{negative sequence}!
        \end{itemize}
    \end{block}
\end{frame}


% ==============================================================================
% SECTION 6: KRON REDUCTION
% ==============================================================================
\section{Kron Reduction (Network Admittance Reduction)}

\begin{frame}[allowframebreaks]{Kron Reduction: Motivation}
    In power system dynamic studies and transient stability simulations:
    \begin{itemize}
        \item Bulk power grids contain tens of thousands of transmission and distribution buses.
        \item However, dynamic states (differential equations) reside almost exclusively in \textbf{synchronous machines} and \textbf{converter-interfaced generation}.
        \item Most intermediate transmission and substation buses have \alert{\textbf{zero current injection}} ($I_{\text{bus}} = 0$).
        \item Passive loads can often be converted to equivalent constant impedances $Y_L = \frac{P_L - j Q_L}{|V|^2}$ at nominal voltage.
    \end{itemize}

    \textbf{The Question Gabriel Kron Answered (1939):}
    Can we algebraically eliminate all non-generator / passive buses while preserving the \alert{\textbf{exact multiport electrical relationships}} between generator internal voltages?
\end{frame}

\begin{frame}[allowframebreaks]{Mathematical Derivation: The Schur Complement}
    Partition the nodal admittance equations $\mathbf{I} = \mathbf{Y}_{\text{bus}} \mathbf{V}$ into:
    \begin{itemize}
        \item Subset $1$ (or $G$): \textbf{Kept buses} (e.g., generator internal EMF nodes).
        \item Subset $2$ (or $L$): \textbf{Eliminated buses} (zero-injection nodes, $\mathbf{I}_2 = \mathbf{0}$).
    \end{itemize}

    \begin{equation}
        \begin{bmatrix} \mathbf{I}_1 \\ \mathbf{0} \end{bmatrix}
        =
        \begin{bmatrix}
            \mathbf{Y}_{11} & \mathbf{Y}_{12} \\
            \mathbf{Y}_{21} & \mathbf{Y}_{22}
        \end{bmatrix}
        \begin{bmatrix} \mathbf{V}_1 \\ \mathbf{V}_2 \end{bmatrix}
    \end{equation}

    From the second block row:
    \begin{equation}
        \mathbf{0} = \mathbf{Y}_{21} \mathbf{V}_1 + \mathbf{Y}_{22} \mathbf{V}_2
    \end{equation}
    Assuming $\mathbf{Y}_{22}$ is invertible (requires at least one shunt path to ground in the eliminated network):
    \begin{equation}
        \mathbf{V}_2 = -\mathbf{Y}_{22}^{-1} \mathbf{Y}_{21} \mathbf{V}_1
    \end{equation}

    Substituting $\mathbf{V}_2$ into the first block row:
    \begin{equation}
        \mathbf{I}_1 = \mathbf{Y}_{11} \mathbf{V}_1 + \mathbf{Y}_{12} \left(-\mathbf{Y}_{22}^{-1} \mathbf{Y}_{21} \mathbf{V}_1\right)
    \end{equation}
    Factoring out $\mathbf{V}_1$:
    \begin{equation}
        \mathbf{I}_1 = \underbrace{\left(\mathbf{Y}_{11} - \mathbf{Y}_{12} \mathbf{Y}_{22}^{-1} \mathbf{Y}_{21}\right)}_{\mathbf{Y}_{\text{Kron}}} \mathbf{V}_1
    \end{equation}

    \begin{tcolorbox}[colback=blue!5,colframe=belgianblue,title=Kron Reduced Admittance Matrix]
        \begin{equation}
            \mathbf{Y}_{\text{Kron}} = \mathbf{Y}_{11} - \mathbf{Y}_{12} \mathbf{Y}_{22}^{-1} \mathbf{Y}_{21}
        \end{equation}
        This is mathematically the \textbf{Schur complement} of block $\mathbf{Y}_{22}$ in $\mathbf{Y}_{\text{bus}}$.
    \end{tcolorbox}
\end{frame}

\begin{frame}[allowframebreaks]{Properties and Applications of Kron Reduction}
    \textbf{1. Exact Multiport Equivalence:}
    \begin{itemize}
        \item The terminal behavior seen from the kept buses is \textbf{100\% mathematically exact}.
        \item No dynamic or steady-state approximation is made (given linear load impedances).
    \end{itemize}

    \textbf{2. Graph-Theoretic Effect: Fill-In:}
    \begin{itemize}
        \item While original power grid admittance matrices are extremely \textbf{sparse} ($>99\%$ zeros), Kron reduction creates equivalent direct transfer admittances between all buses connected through the eliminated network.
        \item $\mathbf{Y}_{\text{Kron}}$ becomes a \alert{\textbf{dense}} matrix!
        \item This corresponds to generalized Star-Delta ($\text{Y}\text{-}\Delta$) transformations.
    \end{itemize}

    \textbf{3. Classical Transient Stability Simulation:}
    \begin{itemize}
        \item In electromechanical stability (the swing equation):
        \begin{equation}
            M_i \frac{d^2\delta_i}{dt^2} = P_{m,i} - P_{e,i}
        \end{equation}
        \item Electrical power output of generator $i$ can be written in closed-form:
        \begin{equation}
            P_{e,i} = E_i^2 G_{ii} + \sum_{j \neq i} E_i E_j \left(G_{ij} \cos(\delta_i - \delta_j) + B_{ij} \sin(\delta_i - \delta_j)\right)
        \end{equation}
        where $Y_{ij} = G_{ij} + j B_{ij}$ are the elements of the \textbf{Kron-reduced admittance matrix}!
    \end{itemize}
\end{frame}

\begin{frame}[allowframebreaks]{Example: 3-Bus System Reduced to 2 Generator Buses}
    Consider two generators (Buses 1 and 2) feeding a common load (Bus 3) with line impedances $z_{13} = j 0.1$, $z_{23} = j 0.2$, and a load admittance $y_{\text{load}} = 2 - j 1$ at Bus 3:
    \begin{center}
        \begin{tikzpicture}[scale=1.0]
            \node[draw, circle, fill=headerbg] (G1) at (0,0) {\scriptsize Gen 1};
            \node[draw, circle, fill=headerbg] (G2) at (4,0) {\scriptsize Gen 2};
            \node[draw, rectangle, fill=lightgraybg] (B3) at (2, -1.3) {\scriptsize Bus 3 (Load)};
            
            \draw[thick] (G1) -- (B3) node[midway, below left] {\tiny $z_{13} = j0.1$};
            \draw[thick] (G2) -- (B3) node[midway, below right] {\tiny $z_{23} = j0.2$};
            \draw[->, thick] (B3) -- (2, -2.0) node[below] {\tiny $y_{\text{load}} = 2 - j1$};
        \end{tikzpicture}
    \end{center}

    Admittances:
    $y_{13} = \frac{1}{j 0.1} = -j 10$, $y_{23} = \frac{1}{j 0.2} = -j 5$.
    The $3 \times 3$ admittance matrix is:
    \begin{equation}
        \mathbf{Y}_{\text{bus}} = \begin{bmatrix}
            -j 10 & 0 & j 10 \\
            0 & -j 5 & j 5 \\
            j 10 & j 5 & (2 - j 16)
        \end{bmatrix}
    \end{equation}

    Eliminating Bus 3:
    $Y_{11} = \begin{bmatrix} -j10 & 0 \\ 0 & -j5 \end{bmatrix}$,
    $Y_{12} = \begin{bmatrix} j10 \\ j5 \end{bmatrix}$,
    $Y_{21} = \begin{bmatrix} j10 & j5 \end{bmatrix}$,
    $Y_{22} = [2 - j16]$.

    \begin{align}
        \mathbf{Y}_{\text{Kron}} &= Y_{11} - Y_{12} Y_{22}^{-1} Y_{21} \\
        &= \begin{bmatrix} -j10 & 0 \\ 0 & -j5 \end{bmatrix} - \frac{1}{2 - j16} \begin{bmatrix} -100 & -50 \\ -50 & -25 \end{bmatrix} \\
        &= \mathbf{\begin{bmatrix} 0.769 - j 3.846 & 0.385 + j 3.077 \\ 0.385 + j 3.077 & 0.192 - j 3.462 \end{bmatrix}}
    \end{align}
    \textbf{Insight:} Gen 1 and Gen 2 are now directly coupled by the transfer admittance $Y_{12}^{\text{Kron}} = 0.385 + j 3.077$. Bus 3 has been completely absorbed!
\end{frame}


% ==============================================================================
% SECTION 7: DISCRETE FOURIER TRANSFORM
% ==============================================================================
\section{Discrete Fourier Transform (DFT) for Phasor Estimation}

\begin{frame}[allowframebreaks]{The Phasor Estimation Problem}
    In the real world:
    \begin{itemize}
        \item Physical voltage and current waveforms are continuous analog signals $x(t)$.
        \item Digital relays, disturbance fault recorders (DFRs), and Phasor Measurement Units (PMUs) sample the waveforms using Analog-to-Digital Converters (ADCs) at sampling frequency $f_s = N \cdot f_0$ (e.g., $N=64$ samples per cycle).
        \item Power system analysis algorithms (symmetrical components, distance protection, power flow, state estimation) require \alert{\textbf{complex phasors}} $\bar{X} = X_{\text{rms}} e^{j\phi}$.
    \end{itemize}

    \textbf{How do we transform discrete sampled time values into accurate fundamental frequency phasors?}
    $\implies$ The \textbf{Full-Cycle Discrete Fourier Transform (DFT)} filter.
\end{frame}

\begin{frame}[allowframebreaks]{The Full-Cycle DFT Filter Formulation}
    Let $x[k] = x(k \Delta t)$ be the sequence of sampled values.
    For a window of $N$ samples over one fundamental cycle:
    \begin{equation}
        \bar{X} = \frac{\sqrt{2}}{N} \sum_{k=0}^{N-1} x[k] e^{-j \frac{2\pi k}{N}}
        = \frac{\sqrt{2}}{N} \sum_{k=0}^{N-1} x[k] \left[\cos\left(\frac{2\pi k}{N}\right) - j \sin\left(\frac{2\pi k}{N}\right)\right]
    \end{equation}

    \begin{columns}
        \begin{column}{0.48\textwidth}
            \begin{block}{Real Part (Cosine Filter)}
                \begin{equation*}
                    X_{\text{Re}} = \frac{\sqrt{2}}{N} \sum_{k=0}^{N-1} x[k] \cos\left(\frac{2\pi k}{N}\right)
                \end{equation*}
            \end{block}
        \end{column}
        \begin{column}{0.48\textwidth}
            \begin{block}{Imaginary Part (Sine Filter)}
                \begin{equation*}
                    X_{\text{Im}} = -\frac{\sqrt{2}}{N} \sum_{k=0}^{N-1} x[k] \sin\left(\frac{2\pi k}{N}\right)
                \end{equation*}
            \end{block}
        \end{column}
    \end{columns}

    \textbf{Phasor Magnitude and Phase:}
    \begin{equation}
        X_{\text{rms}} = \sqrt{X_{\text{Re}}^2 + X_{\text{Im}}^2}, \qquad \phi = \operatorname{atan2}(X_{\text{Im}}, X_{\text{Re}})
    \end{equation}
\end{frame}

\begin{frame}[allowframebreaks]{Filtering Properties of the Full-Cycle DFT}
    \begin{block}{Complete Harmonic Rejection}
        Because $\sum_{k=0}^{N-1} e^{-j \frac{2\pi k m}{N}} = 0$ for any integer $m \neq 0$:
        \begin{itemize}
            \item \textbf{DC offset rejection:} A constant DC component evaluates to zero.
            \item \textbf{Integer harmonics:} 2nd (100 Hz), 3rd (150 Hz), 5th (250 Hz), etc., are \alert{\textbf{completely eliminated}}!
        \end{itemize}
    \end{block}

    \begin{alertblock}{Limitations in Practice}
        \begin{itemize}
            \item \textbf{Decaying DC transients:} Fault currents contain exponentially decaying DC ($I_{\text{dc}} e^{-t/\tau}$). A standard DFT filter leaks some decaying DC into the estimate $\implies$ requires a Mimic Filter ($1 + \tau \frac{d}{dt}$).
            \item \textbf{Off-nominal frequency:} If grid frequency drifts to 49.8 Hz, spectral leakage occurs $\implies$ requires frequency tracking / interpolation.
        \end{itemize}
    \end{alertblock}
\end{frame}


% ==============================================================================
% SECTION 8: SUMMARY & SYNTHESIS
% ==============================================================================
\section{Summary and Conclusions}

\begin{frame}[shrink=12]{Master Summary of Power System Transforms}
    \begin{table}
        \centering
        \tiny
        \begin{tabularx}{\textwidth}{p{2.3cm} p{2.2cm} p{3.2cm} X}
            \toprule
            \textbf{Transform} & \textbf{Coordinates} & \textbf{Mathematical Nature} & \textbf{Key Engineering Benefit} \\
            \midrule
            \textbf{Fortescue} (1918) & Phasor ($0, 1, 2$) & Eigenvector decomposition of cyclic matrices & Decoupled positive, negative, and zero sequence networks; fault calculation; unbalance protection (46, 51N). \\
            \addlinespace
            \textbf{Clarke} (1938) & Stationary ($\alpha, \beta, 0$) & 3-phase to 2-axis orthogonal projection & Space Vector modulation (SVPWM); instant power $p\text{-}q$; fast fault detection. \\
            \addlinespace
            \textbf{Park} (1929) & Synchronous ($d, q, 0$) & Angle-dependent 2D rotation $\mathbf{R}(\theta)$ & Eliminates rotor position from machine equations; converts AC into DC; decoupled $P\text{-}Q$ control; SRF-PLL. \\
            \addlinespace
            \textbf{Kron Reduction} (1939) & Admittance matrix $\mathbf{Y}_{\text{bus}}$ & Schur complement: $\mathbf{Y}_{11} - \mathbf{Y}_{12}\mathbf{Y}_{22}^{-1}\mathbf{Y}_{21}$ & Eliminates interior passive nodes; essential for transient stability simulations and dynamic equivalencing. \\
            \addlinespace
            \textbf{Full-Cycle DFT} & Sampled time to phasor & Discrete orthogonal trigonometric filter & Extracts RMS and phase angle from ADC samples; rejects DC and harmonics; standard in PMUs and numerical distance relays. \\
            \bottomrule
        \end{tabularx}
        \caption{Master reference table of power system transforms.}
    \end{table}
\end{frame}

\begin{frame}[allowframebreaks]{Key Takeaways}
    \begin{enumerate}
        \item \textbf{Decoupling is the primary motivation:} Real three-phase physical components have mutual electromagnetic couplings. Transforms find the natural eigenspace where equations become independent single-phase systems.
        \item \textbf{Choosing the right domain:}
        \begin{itemize}
            \item For steady-state unbalance and fault currents $\to$ \textbf{Fortescue Symmetrical Components}.
            \item For instantaneous converter switching and trajectory visualization $\to$ \textbf{Clarke Transform}.
            \item For closed-loop inverter control and rotating machinery $\to$ \textbf{Park Transform}.
            \item For network simplification and multi-machine dynamics $\to$ \textbf{Kron Reduction}.
            \item For real-time measurement and digital relaying $\to$ \textbf{Full-Cycle DFT}.
        \end{itemize}
        \item \textbf{Cross-domain connections:}
        Positive sequence in Fortescue corresponds to a counter-clockwise circle in $\alpha\beta$, which collapses into a stationary DC point in the $dq$ frame! Negative sequence creates a 100 Hz ripple in $dq$.
    \end{enumerate}
\end{frame}

\begin{frame}
    \centering
    \Huge Thank You!
    \vfill
    \large Questions and Discussion
    \vfill
    \small Python companion code and interactive notebook available in:
    \texttt{Lectures/transforms/transforms.py} and \texttt{transforms.ipynb}
\end{frame}
